\documentclass[11pt,a4paper]{article}
\usepackage[utf8]{inputenc}
\usepackage[T1]{fontenc}
\usepackage{amsmath,amsfonts,amssymb}
\usepackage{graphicx}
\usepackage{hyperref}
\usepackage{listings}
\usepackage{color}
\usepackage{booktabs}
\usepackage{float}
\usepackage{geometry}
\geometry{margin=1in}
\definecolor{zenblue}{RGB}{41,121,255}
\hypersetup{colorlinks=true,linkcolor=zenblue,urlcolor=zenblue,citecolor=zenblue}

\title{\textbf{Zen-Guard-Stream: A Streaming Safety Classifier\\
Fine-Tuned from Qwen2.5-3B}\\[0.5em]
\large Technical Whitepaper v2025.05}
\author{Zach Kelling \\ Zen LM Research Team\\
\texttt{research@zenlm.org}\\
\href{https://papers.zenlm.org}{papers.zenlm.org}}
\date{May 2025}

\begin{document}
\maketitle

\begin{abstract}
Zen-Guard-Stream is a streaming safety classifier built by fine-tuning Alibaba's
\textbf{Qwen2.5-3B} base model~\cite{qwen25}. It is \emph{not} a from-scratch model and uses
no bespoke architecture: the base is \texttt{Qwen/Qwen2.5-3B}, a dense decoder-only
transformer (\texttt{Qwen2ForCausalLM}; 3.09B parameters, 36 layers, hidden size 2048, GQA
with 16 query / 2 key--value heads, vocab 151{,}936, 32K native context). The intended use is
to run alongside a streaming generator and judge the safety of the generation \emph{trajectory}
so far, so that an unsafe continuation can be intercepted before a complete harmful output is
delivered --- something post-hoc filtering cannot do without buffering. \textbf{License
caveat:} unlike most Qwen2.5 sizes, Qwen2.5-3B is released under the \emph{Qwen Research
License} (non-commercial use only), \emph{not} Apache-2.0; any commercial deployment of a
Qwen2.5-3B derivative requires a separate license from Alibaba. We report no safety benchmark
numbers: the upstream base is a general LLM with no published safety metrics, and we have not
run a rigorous streaming-safety evaluation; the figures (e.g. ``ToxiGen 98.7\%,'' ``1.7\,ms
overhead,'' ``0.8\% FPR'') and the ``1.5B-parameter'' size in earlier revisions were
fabricated and have been removed/corrected.
\end{abstract}

\tableofcontents
\newpage

\section{Introduction}

The deployment of streaming generative AI creates a fundamental safety tension: streaming delivers tokens to users immediately as they are generated, enabling low-latency conversational experiences, but this means harmful content reaches users before a post-hoc safety filter can intervene.

Traditional mitigation approaches are inadequate for streaming contexts:

\begin{enumerate}
  \item \textbf{Post-hoc filtering}: Buffer the complete output, apply safety filter, deliver (or withhold) the result. This works for batch generation but adds 100--500ms latency in streaming contexts, negating the latency advantage of streaming.
  \item \textbf{Prompt-level filtering}: Classify the input prompt before generation begins. This misses unsafe content that emerges mid-generation in response to apparently benign prompts (indirect jailbreaks, gradual elicitation).
  \item \textbf{Sampling-time constraints}: Apply logit masks to suppress individual tokens associated with harmful vocabulary. This misses multi-token harmful patterns and is trivially circumvented by paraphrasing.
\end{enumerate}

Zen-Guard-Stream operates at the token generation level with access to the full generation context, enabling detection of harmful trajectories through multi-token patterns while maintaining streaming delivery of safe content.

\subsection{Model Overview}

\begin{table}[H]
\centering
\caption{Zen-Guard-Stream specification. Base-model facts are those of the upstream
Qwen2.5-3B~\cite{qwen25}; Zen-Guard-Stream is a fine-tune of it for trajectory-level safety
scoring.}
\begin{tabular}{ll}
\toprule
\textbf{Parameter} & \textbf{Value} \\
\midrule
Base model & Qwen2.5-3B (Alibaba) \\
License & Qwen Research License (non-commercial only) --- \emph{not} Apache-2.0 \\
Architecture & Dense decoder-only transformer (\texttt{Qwen2ForCausalLM}) \\
Total Parameters & 3.09B \\
Layers / hidden size & 36 / 2048 \\
Attention heads (Q / KV, GQA) & 16 / 2 \\
Vocabulary & 151{,}936 \\
Native context length & 32{,}768 (32K) \\
Integration mode & Trajectory scoring alongside a streaming generator \\
Version & v2025.05 \\
\bottomrule
\end{tabular}
\end{table}

Safety benchmark accuracies, false-positive rates, and per-step overheads are deliberately
omitted (see abstract); the values in earlier revisions were not measured.

\section{Architecture}

\subsection{Streaming Safety Formulation}

Zen-Guard-Stream models streaming safety as a sequential decision problem. At each generation step $t$, given the prompt $x$ and the generated sequence so far $y_{<t}$, the model predicts the safety of continuing generation along the current trajectory:

\begin{equation}
  s_t = f_\theta(x, y_{<t}) \in \{0, 1\}
\end{equation}

where $s_t = 1$ indicates a safe trajectory and $s_t = 0$ triggers an intervention. The model does not classify individual tokens in isolation; it classifies the generation \emph{trajectory} up to and including position $t$, which enables detection of multi-token harmful patterns.

\subsection{Causal Classifier Architecture}

Zen-Guard-Stream is the Qwen2.5-3B causal transformer~\cite{qwen25} with a safety
classification head, run alongside the main generator:

\begin{itemize}
  \item 3.09B parameters in a 36-layer causal transformer (hidden size 2048).
  \item Grouped-query attention (16 query / 2 key--value heads), vocabulary 151{,}936.
  \item A safety head on the final-layer hidden state: $\hat{s}_t = \sigma(w \cdot h_t^{(36)})$,
        i.e. a binary safe/unsafe probability for the trajectory through step $t$.
  \item Intended to run on a separate CUDA stream from the main generator so the two overlap.
\end{itemize}

Running the classifier concurrently with the main generator is a design choice to limit the
marginal latency impact; we do not claim a specific per-step overhead, as it depends on the
hardware, the generator size, and how much computation actually overlaps. The earlier
revision's ``1.5B, 24-layer, 100K-vocabulary, 4-KV-head'' description did not match the
deployed base and has been corrected to the actual Qwen2.5-3B configuration.

\subsection{Intervention Protocol}

When the classifier triggers ($\hat{s}_t < \tau_{\text{safe}} = 0.15$), the streaming intervention protocol executes:

\begin{enumerate}
  \item \textbf{Immediate halt}: The unsafe token is suppressed; streaming delivery pauses.
  \item \textbf{Context assessment}: The main model's logit distribution is inspected; if a safe alternative completion is available with $p > 0.3$, it is substituted and streaming resumes.
  \item \textbf{Graceful degradation}: If no safe continuation is available, a stop message is inserted and the session ends gracefully with an explanation.
  \item \textbf{Logging}: The full context at time of intervention is logged for safety audit review.
\end{enumerate}

The substitution mechanism (step 2) is novel: rather than simply halting generation, Zen-Guard-Stream attempts to steer the generation onto a safe trajectory, improving user experience in borderline cases where the main model has accessible safe alternatives.

\subsection{Context Window}

The classifier sees the recent generation context (prompt plus generated tokens) within the
base model's 32K native context window~\cite{qwen25}, which is more than enough to capture
multi-turn elicitation patterns in a typical conversation. The ``learned context compression
module'' producing a $\text{Compress}_\phi(\cdot)$ summary embedding, described in earlier
revisions, did not exist in the deployed model and has been removed; the model simply uses its
native context window.

\section{Training Methodology}

\subsection{Approach}

Starting from the Qwen2.5-3B base~\cite{qwen25}, Zen-Guard-Stream is fine-tuned to predict a
binary safe/unsafe label for the generation trajectory. The intended objective weights false
negatives (missing unsafe trajectories) more heavily than false positives (blocking safe
creative content),

\begin{equation}
  \mathcal{L} = w_{\text{FN}} \mathcal{L}_{\text{FN}} + w_{\text{FP}} \mathcal{L}_{\text{FP}} + \lambda \mathcal{L}_{\text{calibration}},
\end{equation}

so that recall on genuinely harmful trajectories is prioritized. Training data is constructed
as generation \emph{trajectories} (safe and unsafe) rather than isolated content items, to
match the streaming deployment setting.

We do not publish dataset sizes or composition: the specific figures in earlier revisions (a
500M-trajectory corpus with per-source percentages, and a five-round ``continuous adversarial
training'' curriculum) were fabricated and did not describe a real training run. The honest
statements are that the base is Alibaba's Qwen2.5-3B and that Zen adds a safety fine-tune; we
make no quantitative accuracy claim without a rigorous, reproducible evaluation.

\section{Evaluation}

We intentionally report no benchmark numbers. The upstream Qwen2.5-3B is a general-purpose
LLM with no published safety metrics~\cite{qwen25}, and we have not run a rigorous,
reproducible streaming-safety evaluation of the Zen fine-tune. The classification-accuracy,
false-positive-rate, streaming-overhead, and multi-turn-jailbreak tables in earlier revisions
(e.g. ToxiGen 98.7\%, 0.8\% FPR, 1.7\,ms/step overhead, 90.3\% multi-turn detection) were
fabricated --- they did not come from any measured evaluation --- and have been removed rather
than replaced with invented numbers.

The defensible \emph{qualitative} argument for this design remains: a trajectory-level
classifier inside the generation loop can, in principle, intervene before a complete harmful
output is delivered, which post-hoc filtering cannot do without buffering, and it can use
multi-token context that a per-token logit mask cannot. Whether a given fine-tune actually
achieves high recall at an acceptable false-positive rate and overhead is an empirical
question that must be answered on the operator's own workload (Section~\ref{sec:eval-rec}).

\subsection{Recommended Evaluation Before Deployment}
\label{sec:eval-rec}

Before relying on Zen-Guard-Stream, operators should measure, on their own traffic: detection
recall on unsafe trajectories (including multi-turn elicitation), false-positive rate on
benign creative/medical/security content, and the actual per-step latency overhead on their
hardware and generator. We provide the model and the integration hook; the validation is the
operator's responsibility.

\section{Integration Guide}

\subsection{Hook Interface}

Zen-Guard-Stream is integrated via a pre-token hook in the generation loop:

\begin{lstlisting}[language=Python, caption=Integration Example]
from zen_guard import StreamGuard

guard = StreamGuard.load("zen-guard-stream-v2025.05")

def generation_hook(context: list[int], next_token: int) -> int:
    safety_score = guard.score(context, next_token)
    if safety_score < 0.15:
        # Unsafe trajectory detected
        safe_alt = guard.get_safe_alternative(context)
        if safe_alt is not None:
            return safe_alt  # Substitute safe token
        else:
            return guard.STOP_TOKEN  # Graceful halt
    return next_token  # Pass through safe token
\end{lstlisting}

\subsection{Deployment Topologies}

Three topologies are possible; the right choice and its actual overhead depend on the
operator's hardware and generator, which is why we give no overhead numbers (the figures in
earlier revisions were not measured):

\begin{table}[H]
\centering
\caption{Deployment configurations (qualitative).}
\begin{tabular}{lll}
\toprule
\textbf{Config} & \textbf{Hardware} & \textbf{Trade-off} \\
\midrule
Co-located (recommended) & Same GPU as generator & Lowest added latency; shares GPU \\
Sidecar (separate GPU) & Dedicated accelerator & Decoupled scaling; network hop \\
CPU sidecar & CPU & Cheapest; highest latency, limited concurrency \\
\bottomrule
\end{tabular}
\end{table}

\section{Related Work}

Streaming content safety has been addressed through output filtering \cite{perspective}, circuit breaker mechanisms \cite{circuitbreaker}, and classifier guidance \cite{classifierguidance}. Token-level safety has been studied through vocabulary constraints \cite{vocabconstraint} and watermarking \cite{watermark}. Zen-Guard-Stream applies trajectory-level classification within the generation loop, combining the timeliness of in-loop intervention with multi-token context; the model itself is a fine-tune of Qwen2.5-3B~\cite{qwen25}.

\section{Conclusion}

Zen-Guard-Stream is a streaming safety classifier built by fine-tuning Alibaba's Qwen2.5-3B base model~\cite{qwen25}; it is not a from-scratch model and is 3.09B parameters, not the ``1.5B'' of earlier revisions. Its premise is that a trajectory-level classifier inside the generation loop can intercept an unsafe continuation before a complete harmful output is delivered, which post-hoc filtering cannot do without buffering. We deliberately make no benchmark or latency claims: the upstream base has no published safety metrics, and we have not run a rigorous evaluation, so the inflated accuracy/FPR/overhead figures of earlier revisions have been removed as fabrications. \textbf{Adopters must also heed the license:} Qwen2.5-3B is under the non-commercial Qwen Research License, so commercial use of a derivative requires a separate license from Alibaba. Operators should validate the model on their own traffic before relying on it.

\section*{Attribution and License}

The base weights, training, and capabilities are Alibaba's Qwen2.5-3B (\texttt{Qwen/Qwen2.5-3B}); Zen contributes a safety fine-tune and packaging. Unlike most Qwen2.5 sizes, Qwen2.5-3B is licensed under the \emph{Qwen Research License} (non-commercial only), not Apache-2.0. We thank the Qwen team for releasing the base model.

\begin{thebibliography}{9}
\bibitem{qwen25} Qwen Team, Alibaba (2024). Qwen2.5 Technical Report. arXiv:2412.15115. Base model: \texttt{Qwen/Qwen2.5-3B} (Qwen Research License, non-commercial).
\bibitem{perspective} Lees, A. et al. (2022). A New Generation of Perspective API. arXiv:2202.11176.
\bibitem{circuitbreaker} Zou, A. et al. (2024). Improving Alignment and Robustness with Circuit Breakers. arXiv:2406.04313.
\bibitem{classifierguidance} Dhariwal, P. \& Nichol, A. (2021). Diffusion Models Beat GANs on Image Synthesis. NeurIPS 2021.
\bibitem{vocabconstraint} Krause, B. et al. (2021). GeDi: Generative Discriminator Guided Sequence Generation. EMNLP 2021.
\bibitem{watermark} Kirchenbauer, J. et al. (2023). A Watermark for Large Language Models. ICML 2023.
\end{thebibliography}

\end{document}
